\section{Data}\label{sec:data}

We begin by describing the data used in this chapter. There is no existing dataset that includes multiple views of a single document.\footnote{As of November 2024.} Instead, we evaluate our unified method, described in \S\ref{sec:3-chapter-method}, on 3 existing datasets: DOC2PPT, LongSumm, and Paper-Poster \cite{Fu2021DOC2PPTAP,chandrasekaran2020overview,Qiang2016LearningTG}. These datasets are chosen because they each involve generating a different view of a document. The three datasets and their associated generation tasks are described below.

\paragraph{Slide Generation.} We use the DOC2PPT dataset~\cite{Fu2021DOC2PPTAP}, which contains 5.8K scientific papers in Computer Science and their respective slide decks. As \textcite{Fu2021DOC2PPTAP} do not release data splits or code, we randomly sample 1K examples from this dataset for evaluation. The slides are provided as an image for each slide. We use the Azure OCR tool to extract the text from each slide.\footnote{\scriptsize \url{https://learn.microsoft.com/en-us/azure/ai-services/computer-vision/overview-ocr}}

\paragraph{Blog Generation.} We use the LongSumm dataset~\cite{chandrasekaran2020overview}, which includes blog posts of scientific papers in the Computer Science domain. Since our approach requires no training or supervision, we use the entire training split from Longsumm as our evaluation set. Of the 531 publicly released blog posts in this set, we could only access 505, with the other 26 including broken links or being behind a paywall. Notably, while Longsumm includes a blind test set of 22 papers, this test set only consists of inputs without their reference outputs, thus making it impossible to compute our custom evaluation metric (see \S\ref{sec:eval-metric}). In the interest of completeness and comparison to prior work, we do, however, submit runs from our systems to the leaderboard and report the results of this blind test set in Appendix~\ref{sec:longsumm-blind}.

%Therefore, we use the larger released train split for our evaluation, since our method does not require any training or supervision. This larger set allows for more robust evaluation than the smaller blind test set. . 

\paragraph{Poster Generation.} We use the Paper-Poster dataset~\cite{Qiang2016LearningTG}, which consists of a dataset of 85 papers in Computer Science and Biology, and their respective scientific posters; two examples containing corrupted PDFs are excluded. Although \textcite{Qiang2016LearningTG} release data splits, they do not release code or results for comparison. Given the small size of the dataset, we use it in its entirety for more robust results. While the authors use the source files to extract the text of posters for evaluation, they only release the PDF formats. To preprocess the reference posters, we found that automatic tools to extract text from documents did not handle the visual layout of posters well, so we manually extracted the text of the posters in this dataset. Note that this process was only done to obtain evaluation scores, and that our unsupervised generation method is capable of creating target documents without the need for reference data.%, or this dataset-specific manual extraction step.
%\sjauhar{You manually transcribed the text from the posters?? I hadn't realized. That is both amazing (for obvious reasons), and concerning (if the reviewers say that this will not scale). If there are assumptions that need to be made for your method to work, those need to be listed out clearly in the Method section.}\isabel{I had to manually extract the poster text for the reference data only, it's not necessary for the generation method, since we automatically extract that with azure. Our evaluation metric does assume you have a reference document, but that's true for all of our comparison metrics.}

% \sjauhar{This contradicts what you say about the poster dataset}\isabel{I manually extracted the text of the posters, the papers were automatically extracted. I edited the above paragraph to hopefully make that more clear}

For all 3 datasets, we use the Azure Document Layout tool to extract the text of the input papers.\footnote{\scriptsize \url{https://learn.microsoft.com/en-us/azure/ai-services/document-intelligence/concept-layout}}





